\chapter{Diseases of the Nose and Sinuses}
\label{ch:nose}
The nose serves as the primary airway for respiration, the organ of olfaction, and a conditioning system that warms, humidifies, and filters inspired air. The paranasal sinuses---maxillary, ethmoid, frontal, and sphenoid---are air-filled cavities lined by respiratory mucosa that drain into the nasal cavity. Diseases of the nose and sinuses are extremely prevalent and include allergic, infectious, inflammatory, structural, and neoplastic conditions.

%-------------------------------------------------------------
\section{Allergic and Inflammatory Conditions}
%-------------------------------------------------------------

%-------------------------------------------------------------

%-------------------------------------------------------------

%-------------------------------------------------------------

%-------------------------------------------------------------

%-------------------------------------------------------------

\label{sec:Allergic and Inflammatory Conditions}
%-------------------------------------------------------------%-------------------------------------------------------------

\subsection{Allergic Rhinitis}
\label{sec:Allergic Rhinitis}
Allergic rhinitis is an IgE-mediated inflammatory condition of the nasal mucosa that affects 10--30\% of the global population. The condition results from exposure to airborne allergens (pollens, dust mites, animal dander, molds) triggering IgE-mediated mast cell degranulation. Seasonal allergic rhinitis (hay fever) correlates with pollen seasons, while perennial allergic rhinitis is year-round. Clinical features include sneezing, nasal congestion, rhinorrhea (clear, watery discharge), nasal pruritus, and nasal obstruction; ocular symptoms (itching, tearing, conjunctival injection) are common. Diagnosis is clinical, supported by skin prick testing or specific serum IgE levels. Management includes allergen avoidance, intranasal corticosteroids (first-line for moderate-to-severe disease), second-generation oral antihistamines, leukotriene receptor antagonists, and allergen immunotherapy for refractory or severe cases. Prognosis is favorable with appropriate treatment, though the condition is chronic.

\subsection{Chronic Rhinosinusitis}
\label{sec:Chronic Rhinosinusitis}
Chronic rhinosinusitis is a symptomatic inflammation of the paranasal sinuses persisting for 12 weeks or longer, classified into CRS with nasal polyps (CRSwNP) and CRS without nasal polyps (CRSsNP). The condition results from persistent mucosal inflammation, often with bacterial biofilms; CRSwNP involves type 2 inflammation with eosinophilic infiltration and IL-5/IL-13 pathways. Clinical features include nasal obstruction, facial pain or pressure, purulent nasal discharge, and hyposmia. Diagnosis requires endoscopic or CT evidence of sinus inflammation. Management includes nasal saline irrigation, intranasal corticosteroids, and antibiotics for acute exacerbations; functional endoscopic sinus surgery (FESS) is indicated for refractory disease; biologic therapies (dupilumab, omalizumab) have shown efficacy for CRSwNP. Prognosis is variable, with many patients requiring long-term management.

\subsection{Nasal Polyps}
\label{sec:Nasal Polyps}
Nasal polyps are soft, pale, edematous masses of inflamed sinonasal mucosa that most commonly arise from the ethmoid sinuses and middle meatus, associated with type 2 inflammation, chronic rhinosinusitis, asthma, aspirin-exacerbated respiratory disease, and cystic fibrosis. The condition results from eosinophilic, IgE-mediated mucosal inflammation. Clinical features include nasal obstruction, anosmia or hyposmia, rhinorrhea, and postnasal drip; large or multiple polyps can cause facial deformity. Diagnosis is by anterior rhinoscopy or nasal endoscopy. Management includes intranasal corticosteroids as first-line therapy, short courses of oral corticosteroids, surgical polypectomy, and dupilumab (anti-IL-4R$\alpha$ monoclonal antibody) for refractory disease. Prognosis is variable, with recurrence being common.

\subsection{Non-Allergic Rhinitis}
\label{sec:Non-Allergic Rhinitis}
Non-allergic rhinitis is a chronic nasal condition that encompasses symptoms without an IgE-mediated mechanism, including vasomotor, drug-induced, hormonal, and gustatory forms. Vasomotor rhinitis results from hypersensitivity of trigeminal nerve endings leading to neurogenic inflammation triggered by environmental factors (temperature changes, strong odors, irritants). Drug-induced rhinitis results from overuse of topical decongestants (rhinitis medicamentosa). Clinical features include clear rhinorrhea and nasal congestion. Diagnosis is by exclusion of allergic and infectious causes. Management includes intranasal corticosteroids, ipratropium bromide spray for rhinorrhea, and avoidance of triggers. Prognosis is variable depending on the ability to avoid triggers.

%-------------------------------------------------------------
\section{Infectious Conditions}
%-------------------------------------------------------------

%-------------------------------------------------------------

%-------------------------------------------------------------

%-------------------------------------------------------------

%-------------------------------------------------------------

%-------------------------------------------------------------

\label{sec:Infectious Conditions}
%-------------------------------------------------------------%-------------------------------------------------------------

\subsection{Acute Bacterial Sinusitis}
\label{sec:Acute Bacterial Sinusitis}
Acute bacterial sinusitis is a bacterial infection of the paranasal sinuses lasting less than 4 weeks, most commonly following a viral upper respiratory infection; the most frequent pathogens are \textit{Streptococcus pneumoniae}, \textit{Haemophilus influenzae}, and \textit{Moraxella catarrhalis}. The condition results from bacterial superinfection after viral mucosal inflammation impairs sinus drainage. Clinical features include facial pain or pressure, nasal obstruction, purulent rhinorrhea, and anosmia; diagnosis requires persistent symptoms for more than 10 days without improvement, worsening after initial improvement, or severe symptoms with high fever and purulent discharge for 3--4 consecutive days. Diagnosis is clinical. Management includes amoxicillin-clavulanate as first-line therapy, with nasal saline irrigation, intranasal corticosteroids, and decongestants as adjunctive measures. Prognosis is excellent with appropriate antibiotic therapy; complications (orbital, intracranial) are rare but serious.

\subsection{Epistaxis}
\label{sec:Epistaxis}
Epistaxis is a nosebleed that affects approximately 60\% of the population at some point, with most bleeds originating from Kiesselbach's plexus on the anterior nasal septum (anterior epistaxis). The condition results from local trauma or systemic factors such as dry air, digital trauma, allergic rhinitis, nasal foreign bodies, anticoagulant use, hypertension, coagulopathy, and hereditary hemorrhagic telangiectasia. Clinical features include bleeding from the anterior or posterior nasal cavity; posterior bleeds from the sphenopalatine artery are less common but more severe. Diagnosis is based on clinical examination. Anterior epistaxis is managed with direct pressure, chemical cauterization (silver nitrate), or anterior nasal packing; posterior epistaxis requires posterior packing, endoscopic sphenopalatine artery ligation, or embolization. Prognosis is good, with most cases resolving with simple measures.

\subsection{Fungal Sinusitis}
\label{sec:Fungal Sinusitis}
Fungal sinusitis encompasses a spectrum of conditions ranging from non-invasive forms (fungal ball, allergic fungal sinusitis) to life-threatening invasive forms (acute invasive fungal rhinosinusitis). Non-invasive forms result from fungal colonization or allergic reaction; invasive forms occur in immunocompromised patients (diabetes with ketoacidosis, neutropenia, transplant recipients) and are caused by Mucorales or \textit{Aspergillus} species. Clinical features of acute invasive disease include rapidly progressive facial pain, nasal congestion, black necrotic eschar, and extension to the orbit and brain. Diagnosis of invasive disease requires tissue biopsy. Management of invasive disease involves aggressive surgical debridement, systemic antifungal therapy (amphotericin B), and reversal of immunosuppression. Prognosis for invasive disease is poor with high mortality.

%-------------------------------------------------------------
\section{Structural Disorders}
%-------------------------------------------------------------

%-------------------------------------------------------------

%-------------------------------------------------------------

%-------------------------------------------------------------

%-------------------------------------------------------------

%-------------------------------------------------------------

\label{sec:Structural Disorders}
%-------------------------------------------------------------%-------------------------------------------------------------

\subsection{Choanal Atresia}
\label{sec:Choanal Atresia}
Choanal atresia is a congenital obstruction of the posterior nasal airway caused by a bony (most common), membranous, or mixed atresia plate, and may be associated with CHARGE syndrome. The condition results from failure of the nasobuccal membrane to rupture during embryonic development. Bilateral choanal atresia presents at birth with cyclic cyanosis (worsening with feeding, improving with crying) and severe respiratory distress; unilateral choanal atresia may present later with unilateral nasal obstruction and rhinorrhea. Diagnosis is confirmed by inability to pass a nasogastric tube and by CT imaging. Management is surgical transnasal or transpalatal repair, ideally in the neonatal period for bilateral disease. Prognosis is good following surgical repair.

\subsection{Deviated Nasal Septum}
\label{sec:Deviated Nasal Septum}
Nasal septal deviation is a misalignment of the cartilaginous or bony septum from the midline that is present to some degree in most individuals. The condition results from developmental or traumatic factors. Clinical features include unilateral or bilateral nasal obstruction, which may be worsened by mucosal swelling during allergic rhinitis or upper respiratory infections; deviation may cause recurrent epistaxis, facial pain, and contribute to obstructive sleep apnea. Diagnosis is by anterior rhinoscopy and CT imaging. Management involves septoplasty for severe symptomatic deviation, which may be combined with turbinate reduction for concurrent turbinate hypertrophy. Prognosis is excellent following surgical correction.

\subsection{Nasal Fracture}
\label{sec:Nasal Fracture}
Nasal fractures are the most common facial fracture, accounting for approximately 40\% of all facial fractures, due to the prominent position of the nose; they result from direct trauma (assault, sports injuries, falls). The condition results from direct impact causing bony or cartilaginous disruption. Clinical features include epistaxis, nasal pain, swelling, ecchymosis, and deformity; nasal septal hematoma may coexist and requires urgent drainage to prevent septal necrosis and saddle-nose deformity. Diagnosis is clinical, supplemented by CT for complex fractures. Management involves closed reduction within 7--10 days of injury; severe or comminuted fractures may require open reduction and internal fixation. Prognosis is good with timely treatment.

